% !TEX root = ../../../main.tex
% Source: images/LEC07/IMG_9425--IMG_9428. Subtopics within syllabus §8.6.
\subsection{Analytic functions and the multiplicity of a zero}
\label{sec:analysis-i:analytic-functions}
\begin{pendingreview}
\begin{definition}[Real analytic function]
  \label{thm:analysis-i:analytic-function}
  A function is \emph{analytic at $a$} if, on some neighbourhood of $a$,
  it equals a convergent power series
  \[
    f(x)=\sum_{k=0}^{\infty}b_k(x-a)^k.
  \]
  It is analytic on an open interval if this holds at every point.
\end{definition}
A convergent power series can be differentiated term by term at every
point strictly inside its radius of convergence. We state this
power-series differentiation theorem here without proof. Repeated
application gives $b_k=f^{(k)}(a)/k!$, so an analytic function equals
its Taylor series locally. Taylor's theorem alone does not assert this
equality. Its remainder must tend to zero.

Suppose $f(a)=0$ and its local series has a first nonzero coefficient
$b_m$, where $m\ge1$. Factoring gives
\[
  f(x)=(x-a)^m g(x),\qquad
  g(x)=\sum_{j=0}^{\infty}b_{m+j}(x-a)^j,\qquad g(a)=b_m\ne0.
\]
The shifted series converges on the same neighbourhood. Continuity
makes $g$ nonzero sufficiently close to $a$, so $a$ is an isolated zero.
The integer $m$ is its \emph{multiplicity}. If every coefficient
vanishes, the function is zero on a neighbourhood and this argument
does not assign a finite multiplicity.

\begin{example}[A simple zero of an exponential difference]
  The function $e^x$ never vanishes, but
  \[
    e^x-1=x\left(1+\frac{x}{2!}+\frac{x^2}{3!}+\cdots\right).
  \]
  The factor in parentheses is continuous and has value $1$ at zero.
  Thus zero is a root of multiplicity one of $e^x-1$.
\end{example}

% Original prompt in IMG_9425: write out the analytical forms of sine and
% cosine and derive Euler's identity. The working below completes it.
\begin{example}[Trigonometric series and Euler's identity]
  Taylor's theorem gives
  \[
    \cos x=\sum_{n=0}^{\infty}\frac{(-1)^n x^{2n}}{(2n)!},\qquad
    \sin x=\sum_{n=0}^{\infty}\frac{(-1)^n x^{2n+1}}{(2n+1)!}.
  \]
  Indeed, successive derivatives cycle through sine and cosine with
  signs, and have absolute value at most one. The remainder after
  degree $N$ is bounded by $|x|^{N+1}/(N+1)!$, which tends to zero
  for fixed $x$ by the ratio test.

  Define $e^z=\sum_{k=0}^{\infty}z^k/k!$ for $z\in\C$.
  The ratio test applied to absolute values proves convergence.
  Its even and odd terms can therefore be separated. Since
  $i^{2n}=(-1)^n$ and $i^{2n+1}=i(-1)^n$, the two displayed series give
  \[
    e^{ix}=\cos x+i\sin x.
  \]
  In particular $e^{i\pi}+1=0$. No complex differentiation is used.
\end{example}
\end{pendingreview}

\subsection{Uniform polynomial approximation}
\label{sec:analysis-i:polynomial-approximation}
\begin{pendingreview}
Taylor polynomials use derivatives at one point. Uniform polynomial
approximation asks for small error on an entire interval and requires
only continuity, even for a function with corners.
\begin{theorem}[Weierstrass approximation theorem]
  \label{thm:analysis-i:weierstrass-approximation}
  If $f\in C[a,b]$ with $a<b$, then for every $\varepsilon>0$ there
  exists a polynomial $P$ such that
  \[
    \forall x\in[a,b]\quad |f(x)-P(x)|<\varepsilon.
  \]
\end{theorem}
\end{pendingreview}
\historynote{Sergei Bernstein's 1912 paper offers a proof of
Weierstrass's theorem based on probability~\cite{Bernstein1912}.
The weights below are binomial probabilities. They average nearby
sampled values and concentrate around the input as $n$ grows.
This connects uniform approximation to concentration of random
averages. Sergei Bernstein is a different mathematician from Felix
Bernstein of \autoref{thm:analysis-i:cantor-bernstein}.}
\begin{pendingreview}
\begin{proof}[Proof by Bernstein polynomials]
  First suppose $[a,b]=[0,1]$. For $n\ge1$, define
  \[
    w_{n,k}(x)=\binom nk x^k(1-x)^{n-k},\qquad
    B_nf(x)=\sum_{k=0}^n w_{n,k}(x)f(k/n).
  \]
  Polynomial values at endpoints use the usual constant-power convention.
  The binomial theorem gives $w_{n,k}(x)\ge0$ and
  $\sum_{k=0}^n w_{n,k}(x)=1$ on $[0,1]$.
  The identities $k\binom nk=n\binom{n-1}{k-1}$ and
  $k(k-1)\binom nk=n(n-1)\binom{n-2}{k-2}$ give
  \[
    \sum_{k=0}^n\frac{k}{n}w_{n,k}(x)=x,\qquad
    \sum_{k=0}^n\left(\frac{k}{n}-x\right)^2w_{n,k}(x)
       =\frac{x(1-x)}n\le\frac1{4n}.
  \]
  For the second identity, first sum $k(k-1)$ and $k$ to obtain
  $\sum k^2w_{n,k}=n(n-1)x^2+nx$, and expand the square.
  For $n=1$ the $k(k-1)$ terms are all zero, so the same identity holds.

  By the extreme value theorem, $|f|\le M$ for some $M>0$.
  By uniform continuity, choose $\delta>0$ such that
  $|f(u)-f(v)|<\varepsilon/2$ whenever $|u-v|<\delta$.
  Since the weights sum to one,
  \[
    |B_nf(x)-f(x)|
      \le\sum_{k=0}^n w_{n,k}(x)|f(k/n)-f(x)|.
  \]
  Split the sum according to whether $|k/n-x|<\delta$ or
  $|k/n-x|\ge\delta$. The first part is at most $\varepsilon/2$.
  For the second part, use $|f(k/n)-f(x)|\le2M$ and
  $1\le(k/n-x)^2/\delta^2$. Thus, for every $x\in[0,1]$,
  \[
    |B_nf(x)-f(x)|\le\frac{\varepsilon}{2}
      +\frac{2M}{\delta^2}\frac{x(1-x)}n
      \le\frac{\varepsilon}{2}+\frac{M}{2n\delta^2}.
  \]
  Taking $n>M/(\varepsilon\delta^2)$ makes the last expression
  strictly less than $\varepsilon$. The choice is independent of $x$,
  proving uniform convergence.

  For a general interval, apply this construction to
  $g(t)=f(a+(b-a)t)$ on $[0,1]$. Substituting
  $t=(x-a)/(b-a)$ in an approximating polynomial gives a polynomial
  in $x$ with the same error bound on $[a,b]$.
\end{proof}
The two mechanisms in the proof are different. Uniform continuity
controls the error at nearby samples. The second-moment bound controls
the total weight of distant samples. A corner can therefore be
approximated without differentiating it, as in
\autoref{fig:analysis-i:bernstein-approximation}.
\end{pendingreview}
\begin{figure}[htbp]
  \centering
  \input{\analysisroot/figures/interval-styles.tex}
\begin{tikzpicture}[interval diagram,x=4cm,y=2.5cm]
  \draw[->] (-.07,0)--(1.15,0) node[right=5pt] {$x$};
  \draw[->] (0,-.1)--(0,1.2) node[above=5pt] {$y$};
  \node[below left=5pt] at (0,0) {$0$};
  \draw (0,1)--(.5,0)--(1,1);
  \draw[gray,dashed,domain=0:1,samples=65]
    plot (\x,{1-2*\x+2*\x*\x});
  \draw[gray,domain=0:1,samples=65]
    plot (\x,{(1-\x)^8+6*\x*(1-\x)^7+14*\x^2*(1-\x)^6
      +14*\x^3*(1-\x)^5+14*\x^5*(1-\x)^3
      +14*\x^6*(1-\x)^2+6*\x^7*(1-\x)+\x^8});
  \node[above=5pt] at (.5,.5) {$B_2f$};
  \node[above=5pt] at (.5,.2734375) {$B_8f$};
  \node[below=5pt] at (.5,0) {$1/2$};
  \node[below=5pt] at (1,0) {$1$};
  \node[left=5pt] at (0,1) {$1$};
\end{tikzpicture}

  \caption[Bernstein approximation of a corner.]{\pendingreviewtext{
    The function $f(x)=|2x-1|$ (black) and its Bernstein polynomials
    of degrees $2$ and $8$ (grey). The smooth polynomials approach a
    function that is not differentiable at $x=1/2$. Only $[0,1]$ is shown.}}
  \label{fig:analysis-i:bernstein-approximation}
\end{figure}
